% GENERATED FILE. Edit canonical lessons, then run npm run generate:books.
% canonical-source-hash: fnv1a64:d30a6dcc2c1ffc2c
% canonical-lessons: JA-C369-juuyou, JA-C369-daiji, JA-C369-tokubetsu, JA-C369-youkyuu, JA-C369-houshin

\chapter{What Matters}
\label{ch:ja-what-matters}
\hlchaptermodality{369}

\begin{chapteropening}
\textbf{By the end of this chapter:} \emph{I can say what matters most and what we do about it.}

Complete the last lesson of chapter 369: I can say what matters most and what we do about it.
\end{chapteropening}

\section[\texorpdfstring{\ja{じゅうよう} (jūyō)}{じゅうよう (jūyō)}]{\ja{じゅうよう} (jūyō) — essential; important}

\label{lesson:JA-C369-juuyou}



\begin{quote}
Before the new one: say the Japanese for one of a pair, then the Japanese for to be one-sided.
\end{quote}

\subsection*{You'll want to know: \ja{じゅうよう}}
\textbf{\ja{じゅうよう}} — \emph{jūyō} — \textquotedblleft{}essential; important\textquotedblright{}.

\begin{quote}
\textbf{\ja{しけんが} \ja{ちかく} \ja{なったら、べんきょうが} \ja{じゅうようでしょう。}} — \emph{shiken ga chikaku nattara, benkyō ga jūyō deshō} — \textquotedblleft{}When the exam gets close, studying will be what matters.\textquotedblright{}
\end{quote}

A na-adjective: \ja{な} before a noun.

\begin{cousinweb}
\textbf{\ja{じゅうよう}} (\emph{jūyō}) is Sino-Japanese (kango): two Chinese characters read with their on'yomi, readings borrowed from Chinese. They are \emph{jū}, \textquotedblleft{}important\textquotedblright{} (Middle Chinese \emph{drjowngX}) and \emph{yō}, \textquotedblleft{}need\textquotedblright{} (\emph{'jiewH}) (KANJIDIC2; Middle Chinese in Baxter's transcription).
\end{cousinweb}

\begin{grammarlens}[title={What you've built}]
One more word for saying what happens if.
\end{grammarlens}

\subsection*{Your turn}
\begin{itemize}
\raggedright
  \item \emph{Say it:} \emph{jūyō}
  \item \emph{Say it:} \emph{jūyō}, once more
  \item \emph{Say it:} \emph{shiken ga chikaku nattara, benkyō ga jūyō deshō} — When the exam gets close, studying will be what matters.
  \item \emph{Recall:} say the Japanese for one of a pair, then the Japanese for to be one-sided, then say \emph{jūyō} again
\end{itemize}

\begin{tcolorbox}[breakable,colback=teal!4,colframe=teal!35!black,title={Before you move on}]
What does \ja{じゅうよう} mean? (\textquotedblleft{}Essential; important\textquotedblright{}.) Say it once more.
\end{tcolorbox}
\section[\texorpdfstring{\ja{だいじ} (daiji)}{だいじ (daiji)}]{\ja{だいじ} (daiji) — important; precious}

\label{lesson:JA-C369-daiji}



\begin{quote}
Before the new one: say the Japanese for to be one-sided, then the Japanese for essential.
\end{quote}

\subsection*{You'll want to know: \ja{だいじ}}
\textbf{\ja{だいじ}} — \emph{daiji} — \textquotedblleft{}important; precious\textquotedblright{}.

\begin{quote}
\textbf{\ja{けがを} \ja{したら、やすむ} \ja{ことが} \ja{だいじでしょう。}} — \emph{kega o shitara, yasumu koto ga daiji deshō} — \textquotedblleft{}If you get hurt, resting will be the important thing.\textquotedblright{}
\end{quote}

A na-adjective: \ja{な} before a noun.

\begin{cousinweb}
\textbf{\ja{だいじ}} (\emph{daiji}) is Sino-Japanese (kango): two Chinese characters read with their on'yomi, readings borrowed from Chinese. They are \emph{dai}, \textquotedblleft{}large\textquotedblright{} (Middle Chinese \emph{dajH}) and \emph{ji}, \textquotedblleft{}matter\textquotedblright{} (\emph{dzriH}) (KANJIDIC2; Middle Chinese in Baxter's transcription).
\end{cousinweb}

\begin{grammarlens}[title={What you've built}]
One more word for saying what happens if.
\end{grammarlens}

\subsection*{Your turn}
\begin{itemize}
\raggedright
  \item \emph{Say it:} \emph{daiji}
  \item \emph{Say it:} \emph{daiji}, once more
  \item \emph{Say it:} \emph{kega o shitara, yasumu koto ga daiji deshō} — If you get hurt, resting will be the important thing.
  \item \emph{Recall:} say the Japanese for to be one-sided, then the Japanese for essential, then say \emph{daiji} again
\end{itemize}

\begin{tcolorbox}[breakable,colback=teal!4,colframe=teal!35!black,title={Before you move on}]
What does \ja{だいじ} mean? (\textquotedblleft{}Important; precious\textquotedblright{}.) Say it once more.
\end{tcolorbox}
\section[\texorpdfstring{\ja{とくべつ} (tokubetsu)}{とくべつ (tokubetsu)}]{\ja{とくべつ} (tokubetsu) — special}

\label{lesson:JA-C369-tokubetsu}



\begin{quote}
Before the new one: say the Japanese for essential, then the Japanese for important.
\end{quote}

\subsection*{You'll want to know: \ja{とくべつ}}
\textbf{\ja{とくべつ}} — \emph{tokubetsu} — \textquotedblleft{}special\textquotedblright{}.

\begin{quote}
\textbf{\ja{たんじょうびに} \ja{なったら、しょくじは} \ja{とくべつです。}} — \emph{tanjōbi ni nattara, shokuji wa tokubetsu desu} — \textquotedblleft{}On my birthday, the meal is special.\textquotedblright{}
\end{quote}

A na-adjective: \ja{な} before a noun.

\begin{cousinweb}
\textbf{\ja{とくべつ}} (\emph{tokubetsu}) is Sino-Japanese (kango): two Chinese characters read with their on'yomi, readings borrowed from Chinese. They are \emph{toku}, \textquotedblleft{}special\textquotedblright{} (Middle Chinese \emph{dok}) and \emph{betsu}, \textquotedblleft{}separate\textquotedblright{} (\emph{bjet}) (KANJIDIC2; Middle Chinese in Baxter's transcription).
\end{cousinweb}

\begin{grammarlens}[title={What you've built}]
One more word for saying what happens if.
\end{grammarlens}

\subsection*{Your turn}
\begin{itemize}
\raggedright
  \item \emph{Say it:} \emph{tokubetsu}
  \item \emph{Say it:} \emph{tokubetsu}, once more
  \item \emph{Say it:} \emph{tanjōbi ni nattara, shokuji wa tokubetsu desu} — On my birthday, the meal is special.
  \item \emph{Recall:} say the Japanese for essential, then the Japanese for important, then say \emph{tokubetsu} again
\end{itemize}

\begin{tcolorbox}[breakable,colback=teal!4,colframe=teal!35!black,title={Before you move on}]
What does \ja{とくべつ} mean? (\textquotedblleft{}Special\textquotedblright{}.) Say it once more.
\end{tcolorbox}
\section[\texorpdfstring{\ja{ようきゅう} (yōkyū)}{ようきゅう (yōkyū)}]{\ja{ようきゅう} (yōkyū) — a demand; a firm request}

\label{lesson:JA-C369-youkyuu}



\begin{quote}
Before the new one: say the Japanese for important, then the Japanese for special.
\end{quote}

\subsection*{You'll want to know: \ja{ようきゅう}}
\textbf{\ja{ようきゅう}} — \emph{yōkyū} — \textquotedblleft{}a demand; a firm request\textquotedblright{}.

\begin{quote}
\textbf{\ja{ようきゅうが} \ja{あったら、かいしゃに} \ja{言}\ja{って} \ja{ください。}} — \emph{yōkyū ga attara, kaisha ni itte kudasai} — \textquotedblleft{}If you have any demands, please tell the company.\textquotedblright{}
\end{quote}

\begin{cousinweb}
\textbf{\ja{ようきゅう}} (\emph{yōkyū}) is Sino-Japanese (kango): two Chinese characters read with their on'yomi, readings borrowed from Chinese. They are \emph{yō}, \textquotedblleft{}need\textquotedblright{} (Middle Chinese \emph{'jiewH}) and \emph{kyū}, \textquotedblleft{}request\textquotedblright{} (KANJIDIC2; Middle Chinese in Baxter's transcription).
\end{cousinweb}

\begin{grammarlens}[title={What you've built}]
One more word for saying what happens if.
\end{grammarlens}

\subsection*{Your turn}
\begin{itemize}
\raggedright
  \item \emph{Say it:} \emph{yōkyū}
  \item \emph{Say it:} \emph{yōkyū}, once more
  \item \emph{Say it:} \emph{yōkyū ga attara, kaisha ni itte kudasai} — If you have any demands, please tell the company.
  \item \emph{Recall:} say the Japanese for important, then the Japanese for special, then say \emph{yōkyū} again
\end{itemize}

\begin{tcolorbox}[breakable,colback=teal!4,colframe=teal!35!black,title={Before you move on}]
What does \ja{ようきゅう} mean? (\textquotedblleft{}A demand; a firm request\textquotedblright{}.) Say it once more.
\end{tcolorbox}
\section[\texorpdfstring{\ja{ほうしん} (hōshin)}{ほうしん (hōshin)}]{\ja{ほうしん} (hōshin) — a policy; a course of action}

\label{lesson:JA-C369-houshin}



\begin{quote}
Before the new one: say the Japanese for special, then the Japanese for a demand.
\end{quote}

\subsection*{You'll want to know: \ja{ほうしん}}
\textbf{\ja{ほうしん}} — \emph{hōshin} — \textquotedblleft{}a policy; a course of action\textquotedblright{}.

\begin{quote}
\textbf{\ja{ほうしんは、しゃちょうの} \ja{いけんに} \ja{よります。}} — \emph{hōshin wa, shachō no iken ni yorimasu} — \textquotedblleft{}The policy depends on what the company president thinks.\textquotedblright{}
\end{quote}

\begin{cousinweb}
\textbf{\ja{ほうしん}} (\emph{hōshin}) is Sino-Japanese (kango): two Chinese characters read with their on'yomi, readings borrowed from Chinese. They are \emph{hō}, \textquotedblleft{}direction\textquotedblright{} (Middle Chinese \emph{pjang}) and \emph{shin}, \textquotedblleft{}needle\textquotedblright{} (\emph{tsyim}) (KANJIDIC2; Middle Chinese in Baxter's transcription).
\end{cousinweb}

\begin{grammarlens}[title={What you've built}]
One more word for saying what happens if.
\end{grammarlens}

\subsection*{Your turn}
\begin{itemize}
\raggedright
  \item \emph{Say it:} \emph{hōshin}
  \item \emph{Say it:} \emph{hōshin}, once more
  \item \emph{Say it:} \emph{hōshin wa, shachō no iken ni yorimasu} — The policy depends on what the company president thinks.
  \item \emph{Recall:} say the Japanese for special, then the Japanese for a demand, then say \emph{hōshin} again
\end{itemize}

\begin{tcolorbox}[breakable,colback=teal!4,colframe=teal!35!black,title={Before you move on}]
What does \ja{ほうしん} mean? (\textquotedblleft{}A policy; a course of action\textquotedblright{}.) Say it once more.
\end{tcolorbox}
